% Appendix: discretization, time integration, verification.
% Every statement here corresponds to code in fusion_ode_identification/
% imex_solver.py and scripts/verify_solver.py, and to artifacts in
% logs/solver_verification/verification_report.json.

\section{Spatial discretization, time integration, and verification}
\label{app:numerics}

\subsection{Grid and control volumes}

The radial domain is discretized with $N = 65$ nodes on the uniform grid
$\rho_i = i/(N-1)$, $i = 0,\dots,N-1$, with spacing
$\Delta\rho = 1/(N-1) \approx 0.0156$. The scheme is node-centered:
unknowns live at nodes and faces sit midway between them, at
$\rho_{i+1/2} = \rho_i + \Delta\rho/2$. The last node, $\rho_{N-1} = 1$, is
the Dirichlet boundary and is not an unknown, so the solve has $N-1 = 64$
degrees of freedom per time level plus the scalar latent.

Control volumes follow from the face positions. Interior cells span one
grid spacing, while the axis cell spans only half a spacing because it is
bounded on the left by the symmetry axis:
\begin{equation}
v_i = V'(\rho_i)\,\Delta\rho \quad (1 \le i \le N-2),
\qquad
v_0 = V'\!\left(\tfrac{\Delta\rho}{4}\right)\frac{\Delta\rho}{2},
\label{eq:cellvol}
\end{equation}
where the axis value is evaluated at the cell centroid $\rho = \Delta\rho/4$
and is obtained in code by the linear blend
$\tfrac{3}{4}V'_0 + \tfrac{1}{4}V'_1$
(\texttt{imex\_solver.cell\_volumes}).

\paragraph{A defect the verification suite caught.} An earlier
implementation used the face average $\tfrac12(V'_i + V'_{i+1})$ for every
cell volume and gave the axis cell full width $\Delta\rho$. Both choices
are $O(\Delta\rho)$ inconsistent with the face positions actually used in
the flux stencil. The resulting scheme was stable and trained without any
visible symptom, but its solution error \emph{did not decrease under grid
refinement}: measured $L_2$ error on the analytic test of
\S\ref{app:num-verif} was $5.4\times10^{-3}$, $8.1\times10^{-3}$,
$9.7\times10^{-3}$, $1.0\times10^{-2}$ for $N = 17, 33, 65, 129$, i.e.\ an
apparent convergence order of $-0.31$. Only the refinement study exposed
this; no training or validation metric did. After adopting
\eqref{eq:cellvol} the same test converges at first order. We report this
because it illustrates why identified physical parameters should not be
interpreted before a convergence study is run: a discretization can be
consistent-looking, stable, trainable, and still wrong.

\subsection{Flux form and the discrete operator}

With $T_i \equiv \Te(\rho_i, t)$, face fluxes use the two-point gradient
and face-averaged coefficients,
\begin{equation}
F_{i+1/2} = -\,V'_{i+1/2}\,\chi_{i+1/2}\,\frac{T_{i+1} - T_i}{\Delta\rho},
\qquad
V'_{i+1/2} = \tfrac12(V'_i + V'_{i+1}),
\quad
\chi_{i+1/2} = \tfrac12(\chi_i + \chi_{i+1}),
\label{eq:faceflux}
\end{equation}
and the semi-discrete divergence on cell $i$ is the flux difference over
the control volume,
\begin{equation}
(L T)_i = -\,\frac{F_{i+1/2} - F_{i-1/2}}{v_i}.
\label{eq:div}
\end{equation}
Defining $k_{i+1/2} = V'_{i+1/2}\chi_{i+1/2}/\Delta\rho$, \eqref{eq:div}
expands to the tridiagonal stencil
\begin{equation}
(L T)_i = \frac{k_{i-1/2} T_{i-1} - (k_{i-1/2} + k_{i+1/2}) T_i + k_{i+1/2} T_{i+1}}{v_i},
\label{eq:stencil}
\end{equation}
so $L$ has non-positive diagonal and non-negative off-diagonal entries: it
is an M-matrix generator, which is what makes the implicit solve of
\S\ref{app:num-imex} unconditionally solvable by the Thomas algorithm
without pivoting.

\paragraph{Boundary rows.} Two rows differ from \eqref{eq:stencil}.
\emph{Axis ($i=0$)}: the symmetry condition
$\partial_\rho \Te|_{\rho=0} = 0$ is imposed as a zero left-face flux,
$F_{-1/2} \equiv 0$. This is a flux condition, not a ghost-node condition:
no extrapolated value is introduced, and the discrete conservation
statement below therefore holds exactly. The row becomes
$(LT)_0 = k_{1/2}(T_1 - T_0)/v_0$.
\emph{Edge ($i=N-2$)}: the outer neighbour $T_{N-1}$ is the prescribed
boundary trace, so its contribution moves to the right-hand side. The
super-diagonal entry of the last row is zero and the boundary coupling
vector has the single nonzero component
$b^{\mathrm{bc}}_{N-2} = k_{N-3/2}/v_{N-2}$, which multiplies
$T_{\mathrm{edge}}$.

\paragraph{Numerical floors.} Three guards act before assembly, all
recorded here because they alter the operator when triggered:
$V' \leftarrow \max(V', 10^{-6})$;
$\Delta\rho \leftarrow \max(\Delta\rho, 10^{-6}\max\Delta\rho + 10^{-12})$;
and $v_i \leftarrow \max(v_i, \max(10^{-4}\max_j v_j, 10^{-10}))$. On the
uniform grid with the equilibrium $V'$ of this corpus none of these floors
activates; per-shot flags recording activation are written into the loss
diagnostics.

\subsection{IMEX time integration}
\label{app:num-imex}

The semi-discrete system is
\begin{equation}
\dot{\bm{T}} = L(\lat)\,\bm{T} + \bm{b}^{\mathrm{bc}}\,T_{\mathrm{edge}}(t)
             + \bm{s}_\theta(\bm{T}, \nel, \inr, t),
\qquad
\dot\lat = F(\lat, \inr)/1 ,
\label{eq:semidiscrete}
\end{equation}
with $F$ as in the main text. Diffusion is stiff: the spectral radius of
$L$ scales as $\chi/\Delta\rho^2$, which for
$\chi \sim 1\,\mathrm{m^2 s^{-1}}$ and $\Delta\rho \approx 0.0156$ is of
order $4\times10^{3}\,\mathrm{s^{-1}}$, so an explicit method would need
$\Delta t \lesssim 2\times10^{-4}\,$s. We therefore treat $L$ implicitly
and everything else explicitly. Each substep of size $\Delta t$ advances
\begin{align}
\big[I - \thIMEX \Delta t\, L(\lat_n)\big] \bm{T}_{n+1}
&= \bm{T}_n
 + \Delta t\,(1-\thIMEX)\, L(\lat_n)\,\bm{T}_n
 + \Delta t\, \bm{s}_\theta(\bm{T}_n,\cdot)
 + \Delta t\, \thIMEX\, \bm{b}^{\mathrm{bc}} T_{\mathrm{edge}}(t_{n+1}),
\label{eq:imexT}\\
\lat_{n+1} &= \lat_n + \Delta t\, F(\lat_n, \inr_n),
\label{eq:imexZ}
\end{align}
with $\thIMEX = 0.7$. Three details of \eqref{eq:imexT}--\eqref{eq:imexZ}
are easy to state ambiguously and are therefore made explicit:
(i) the operator is frozen at the \emph{current} latent state $\lat_n$, so
the implicit stage is linear in $\bm{T}$ and no nonlinear iteration is
needed --- this is the reason $\chi$ is allowed to depend on $\lat$ and
$\rho$ but never on $\Te$;
(ii) the boundary term is evaluated at the \emph{new} time level, weighted
by $\thIMEX$, consistently with the implicit part;
(iii) the latent uses explicit Euler at the same substep size, evaluated
before the temperature update, so the pair is a first-order Lie splitting.

Within each observation interval $[t_j, t_{j+1}]$ the solver takes a fixed
number $M = 5$ of substeps, $\Delta t = (t_{j+1}-t_j)/M$, and exogenous
inputs are linearly blended across the interval. Fixed $M$ (rather than
adaptive stepping) keeps all loop bounds static, which is what permits
reverse-mode differentiation through the entire rollout under JIT
compilation. Past the valid window of a discharge the state is frozen by
an activity mask so that zero padding cannot destabilize the solve or
contaminate gradients.

\paragraph{Stability and order.} For the scalar test problem
$\dot y = \lambda y$ the update \eqref{eq:imexT} has amplification factor
\begin{equation}
R(s) = \frac{1 + (1-\thIMEX)s}{1 - \thIMEX s}, \qquad s = \lambda \Delta t,
\label{eq:amp}
\end{equation}
which satisfies $|R| \le 1$ for all $\mathrm{Re}\,s \le 0$ when
$\thIMEX \ge 1/2$: the method is A-stable in that range. It is not
L-stable for $\thIMEX < 1$, since
$R(s) \to -(1-\thIMEX)/\thIMEX = -0.43$ as $s \to -\infty$; the stiffest
modes are damped by a factor $0.43$ per substep rather than annihilated.
Backward Euler ($\thIMEX = 1$) would be L-stable but is more dissipative
at moderate $s$; Crank--Nicolson ($\thIMEX = 1/2$) is second order but
$R \to -1$, so stiff modes ring rather than decay. The choice
$\thIMEX = 0.7$ is a compromise, and it makes the scheme
\emph{first order} in time, not second: the temporal error observed in
\S\ref{app:num-verif} is consistent with this. We state this explicitly
because a ``$\theta$-method'' is often assumed to be second order.
Implicit treatment removes the explicit parabolic step-size restriction;
it does not remove the physical stiffness, and the explicit source and
latent terms retain their own accuracy-driven step limits.

\subsection{Verification}
\label{app:num-verif}

All tests below are implemented in \texttt{scripts/verify\_solver.py} and
run on the same code path used for training. Results are archived in
\texttt{logs/solver\_verification/verification\_report.json}.

\paragraph{Analytic reference solution.} With $V' \equiv 1$,
$\chi \equiv \chi_0$, homogeneous Neumann at $\rho=0$ and homogeneous
Dirichlet at $\rho=1$, the function
\begin{equation}
\Te(\rho,t) = A\cos(k\rho)\,e^{-\chi_0 k^2 t}, \qquad k = \pi/2,
\label{eq:analytic}
\end{equation}
solves the continuous problem exactly and satisfies both boundary
conditions. We integrate to $t_{\mathrm{end}} = 0.05\,$s with
$\chi_0 = 1$ and measure the relative $L_2$ error at the final time.
The amplitude is set to $A = 1000\,$eV rather than unity: the solver's
smooth positivity safeguard is exact only to $\sim0.01\,$eV, so a
unit-amplitude test measures the safeguard rather than the discretization.
This is itself a small lesson --- a verification test must be run at the
operating magnitude of the quantity it verifies.

\paragraph{Grid refinement.} With $M = 40$ substeps (temporal error
subdominant):
\begin{center}
\begin{tabular}{lrrrr}
\toprule
$N$ & 17 & 33 & 65 & 129 \\
$\Delta\rho$ & 0.0625 & 0.0312 & 0.0156 & 0.0078 \\
relative $L_2$ error & $3.47\times10^{-3}$ & $1.68\times10^{-3}$ & $8.26\times10^{-4}$ & $4.10\times10^{-4}$ \\
\bottomrule
\end{tabular}
\end{center}
A least-squares fit of $\log\varepsilon$ against $\log\Delta\rho$ gives
observed order $1.03$. The error halves with each refinement, confirming
consistency; the first-order rate reflects the splitting and the
first-order treatment of the boundary/latent coupling rather than the
interior stencil, which is formally second order.

\paragraph{Substep refinement.} At $N = 129$:
\begin{center}
\begin{tabular}{lrrrrr}
\toprule
$M$ & 1 & 2 & 5 & 10 & 20 \\
relative $L_2$ error & $4.72\times10^{-4}$ & $4.40\times10^{-4}$ & $4.21\times10^{-4}$ & $4.15\times10^{-4}$ & $4.12\times10^{-4}$ \\
\bottomrule
\end{tabular}
\end{center}
The temporal error saturates against the spatial floor by $M \approx 5$,
which is why $M = 5$ is used in training: further substeps buy accuracy
that the spatial resolution cannot express, at linear cost.

\paragraph{Operating-point error budget.} At the training configuration
($N = 65$, $M = 5$) the solver's own error on \eqref{eq:analytic} is
$8.3\times10^{-4}$ relative, i.e.\ well below $1\,$eV on a $1000\,$eV
scale. The model--data misfit reported in the main text is
$\sim\resMAE$, two orders of magnitude larger. Numerical error is
therefore not a limiting factor in any comparison we draw.

\paragraph{Discrete conservation.} Summing \eqref{eq:div} against the cell
volumes must telescope to the net boundary flux. On a parabolic test
profile the measured volume-weighted sum
$\sum_i v_i (LT)_i$ equals $-F_{N-3/2}$ to a relative residual of
$0$ (exactly, in double precision, since both sides are the same finite
sum). This verifies that the axis condition contributes no spurious flux
and that the cell volumes \eqref{eq:cellvol} are consistent with the face
positions in \eqref{eq:faceflux} --- the property the earlier
implementation violated.

\paragraph{Positivity.} Starting from a non-negative step profile with a
sharp interior discontinuity and a hot Dirichlet edge, the minimum value
over all times and nodes remains $\ge 0$ (measured $0.0$). The scheme is
not formally positivity-preserving for arbitrary $\Delta t$, but the
M-matrix structure of $I - \thIMEX\Delta t L$ together with the smooth
clamp keeps solutions in the physical range at the step sizes used.

\paragraph{Gradient correctness.} Because parameters are identified by
differentiating through \eqref{eq:imexT}--\eqref{eq:imexZ}, the adjoint
must be verified, not assumed. We compare reverse-mode automatic
derivatives of the per-discharge loss against central finite differences
($\epsilon = 10^{-5}$) on one real training discharge:
\begin{center}
\begin{tabular}{lrrr}
\toprule
parameter & autodiff & central difference & relative error \\
\midrule
$\beta$ (raw)               & $-1.15597\times10^{-1}$ & $-1.15597\times10^{-1}$ & $6.6\times10^{-10}$ \\
$\tau$ (raw)                & $-1.33876\times10^{-1}$ & $-1.33876\times10^{-1}$ & $9.3\times10^{-10}$ \\
drive intercept $\kappa_0$  & $+5.08148\times10^{-1}$ & $+5.08148\times10^{-1}$ & $6.7\times10^{-10}$ \\
$\chi$ edge gap (raw)       & $+1.27057\times10^{-1}$ & $+1.27057\times10^{-1}$ & $1.9\times10^{-10}$ \\
\bottomrule
\end{tabular}
\end{center}
All agree to nine or ten digits. Differentiation is through the
\emph{discrete} solver (a discrete adjoint obtained by automatic
differentiation); no continuous adjoint equation is solved, and we avoid
that terminology deliberately.

\paragraph{Solver robustness.} Across every discharge and every training
run reported here, the fraction of rollouts flagged as failed (non-finite
state or step-count overflow) was zero; the success indicator entering the
loss was $1.00$ at every logged step.
